\begin{textsample}{NEG Dim 2 – 2013\_c5 – Score -2.00 – 2013\_c5/enade\_2013\_farm\_d0638aebea.txt}  \label{ex:f2_neg_064}
CAPÍTULO 5 DISTRIBUIÇÃO DOS CONCEITOS 5.1 PANORAMA NACIONAL DA DISTRIBUIÇÃO DOS CONCEITOS A Tabela 5.1 apresenta a quantidade e distribuição de cursos de Farmácia participantes do ENADE/2013, por faixa de conceito e Grande Região.
A diferença entre os cursos tabulados neste capítulo e no capítulo 2 são os cursos sem conceito, em princípio, aqueles sem alunos concluintes que participassem da prova ou que tivessem tão somente um aluno concluinte13.
Observando -se os dados da Tabela 5.1, nota -se que, dos 309 cursos participantes, 107 ( 34,6 \% ) classificaram -se com conceito 3, o valor modal.
Este foi também o conceito modal nas regiões Nordeste, Sudeste e Sul do Brasil ( entre 30,9 \% na região Sudeste e 48,4 \% na região Sul ), enquanto nas regiões Norte e Centro-Oeste, a moda foi o conceito 2 ( 58,8 \% e 32,4 \% ), respectivamente.
O conceito 2 foi o segundo mais frequente em nível nacional ( 29,8 \%, correspondendo a 92 cursos ) e o conceito 4, o terceiro ( 21,0 \%, correspondendo a 65 cursos ).
Houve, ainda, 14 cursos ( 4,5 \% ) que receberam conceito 5 e outros 20 cursos que receberam conceito 1.
Dos 309 cursos de Farmácia, 11 ficaram sem conceito ( SC ).
A região Norte participou com 17 cursos ou 5,5 \% do total nacional.
Nesta região, o conceito 3 foi atribuído a quatro cursos dos 17 participantes, o que equivale a 23,5 \% do total regional.
Apenas um curso ( 5,9 \% ) ficou sem conceito ( SC ) e outros dois cursos alocados nos conceitos 1 e 4 ( um curso, cada ).
Dez cursos foram avaliados com conceito 2 ( 58,8 \% ), o conceito modal na região.
Nenhum curso foi alocado ao conceito 5. 13 Estes cursos com somente um concluinte não foram considerados no capítulo 2.
A região Nordeste participou com 44 cursos ou 14,2 \% do total nacional.
Nesta região, 16 cursos, 36,4 \% em termos regionais, obtiveram conceito 3, o conceito modal para a região.
Os conceitos 2 e 4 foram atribuídos a, respectivamente, 14 e dez cursos ( 31,8 \% e 22,7 \% ).
Já os conceitos 1 e 5 foram atribuídos a, respectivamente, um e dois cursos ( 2,3 \% e 4,5 \% ).
Nessa região, apenas um dos cursos ( 2,3 \% ) ficou sem conceito ( SC ).
Dos 152 cursos participantes da região Sudeste, 47 ( 30,9 \% ) obtiveram conceito 3, o conceito modal.
O conceito 1 foi atribuído a 14 cursos ( 9,2 \% ) e o conceito 2, a 45 ( 29,6 \% ).
O conceito 4 foi atribuído a 31 cursos ( 20,4 \% ) e sete outros ( 4,6 \% ) receberam o conceito 5.
Apenas oito cursos ficaram sem conceito ( 5,3 \% ).
A região Sul também contou com cursos \textbf{distribuídos} em todas as faixas de conceitos.
A predominância do conceito 3 foi de 48,4 \%, como já comentado anteriormente, correspondentes a 30 dos 62 cursos participantes na região Sul.
O conceito 2 foi atribuído a 12 cursos ( 19,4 \% ) e o conceito 4, a 14 cursos ( 22,6 \% ).
O conceito 1 foi recebido por um único curso ( 1,6 \% ) e o conceito 5, por cinco deles ( 8,1 \% ).
Nenhum dos cursos da região Sul ficou sem conceito ( SC ).
Como já comentado, de todos os 34 cursos participantes na região Centro-Oeste, 11 receberam conceito 2, o conceito modal ( 32,4 \% ).
Nesta região nenhum curso ficou com conceito 5, enquanto apenas um curso ficou sem conceito ( SC ).
Os demais cursos foram avaliados com conceitos 1 ( três cursos, 8,8 \% ), 3 ( dez cursos, 29,4 \% ) e 4 ( nove cursos, 26,5 \% ).
O Gráfico 5.1 apresenta a distribuição do Conceito ENADE segundo Grande Região : poligonais mais à esquerda representam Grandes Regiões com pior distribuição, e poligonais mais à direita, Grandes Regiões com melhores conceitos.
A região Norte ( linha azul ) apresenta os cursos com valores mais baixos, ao passo que as regiões Sul ( linha roxa ) e Nordeste ( linha verde ) apresentam as distribuições com valores maiores.
Já os cursos em IES das regiões Sudeste ( linha vermelha ) e Centro-Oeste ( linha preta ) apresentam uma situação intermediária. 5.2 CONCEITOS POR CATEGORIA ADMINISTRATIVA E POR GRANDE REGIÃO A Tabela 5.2 apresenta a distribuição dos cursos participantes do ENADE/2013 de Farmácia, por Categoria Administrativa, de acordo com os conceitos por eles alcançados, segundo as Grandes Regiões brasileiras.
Dos 309 cursos participantes, 60 ( 19,4 \% ) eram ministrados em Instituições Públicas e 249 ( 80,6 \% ), em Privadas.
De acordo com as informações da Tabela 5.2, em termos nacionais, entre Instituições Públicas e Privadas : das 14 IES que receberam o conceito 5, 11 eram públicas e três privadas.
Dos 60 cursos participantes de IES Públicas, o conceito 4 foi o valor modal, atribuído a 32 cursos ( 53,3 \% ).
Entre os demais cursos de Instituições Públicas participantes, somente um obteve conceito 1 ( 1,7 \% da categoria ), sete receberam conceito 2 ( 11,7 \% ), nove foram alocados ao conceito 3 ( 15,0 \% ) e os 11 restantes receberam o conceito 5 ( 18,3 \% ).
Nesta categoria, nenhum dos cursos ficou sem conceito ( SC ).
Na rede privada, o conceito modal foi 3, com 98 cursos dos 249 da categoria.
Entre os demais cursos participantes, 19 receberam conceito 1 e 85, conceito 2.
O conceito 4 foi atribuído a 33 cursos e o conceito 5 a três cursos.
Nesta Categoria Administrativa, 11 dos cursos ficaram sem conceito ( SC ).
Na análise por região, observa -se que, na região Norte, as Instituições Privadas participaram com 13 cursos ( 76,5 \% do total regional ), dos quais apenas um ficou sem conceito ( SC ).
O conceito modal para as Instituições Privadas na região foi 2, com nove cursos, correspondendo a 69,2 \%.
Seguida do conceito 3, com três cursos, correspondendo a 23,1 \%.
As Instituições Públicas participaram com quatro cursos ( 23,5 \% do total regional ), dos quais nenhum ficou sem conceito ( SC ).
Todos os cursos ficaram alocados igualmente aos conceitos 1, 2, 3 e 4 ( um curso, cada ).
Na região Nordeste, a rede privada concentrou 28 dos 44 cursos participantes, equivalentes a 63,6 \% do total da região, a menor proporção entre todas as regiões.
Nesta Categoria Administrativa na região, \textbf{predominaram} os cursos com conceito 2, 14 cursos correspondendo a 50,0 \% dos cursos oferecidos por IES Privadas no Nordeste.
Apenas um curso recebeu o conceito 1 ( 3,6 \% ), enquanto o conceito 5 não foi atribuído a nenhum curso.
Os demais foram avaliados com conceitos 3 ( 11 cursos correspondendo a 39,3 \% ) e 4 ( também um curso correspondendo a 3,6 \% ).
As Instituições Públicas dessa região participaram com 16 cursos ( 36,4 \% ), dos quais nove obtiveram conceitos 4, o valor modal.
Os demais foram avaliados com os conceitos 3 ( cinco cursos ) e 5 ( dois cursos ).
Nenhum curso ficou sem conceito ( SC ) ou receberam conceitos 1 ou 2 nesta combinação de região e categoria.
Na região Sudeste, a proporção de cursos da rede privada, 87,5 \%, foi mais elevada do que nas demais regiões brasileiras, correspondendo a 133 dos 152 cursos participantes.
Nesta categoria, na região Sudeste, o conceito modal foi 3 ( 46 cursos ).
Nesta combinação de Categoria Administrativa e Grande Região, apenas oito dos cursos não receberam conceito ( SC ).
Os demais foram avaliados com conceito 1 ( 14 cursos ), conceito 2 ( 43 cursos ), conceito 4 ( 19 cursos ) e conceito 5 ( três cursos ).
Entre os 19 cursos em Instituições Públicas na região Sudeste, a categoria modal foi a que recebeu conceito 4 ( 12 cursos ).
Os demais foram avaliados com os conceitos 2 ( dois cursos ), 3 ( um curso ) e 5 ( quatro cursos ).
Na região Sudeste, nesta categoria, nenhum curso ficou sem conceito ( SC ) ou recebeu conceito 1.
As Instituições Privadas concentraram 48 dos 62 cursos participantes da região Sul, 77,4 \% do total regional.
Desses, 28 ficaram com conceito 3, o conceito modal.
Nesta combinação de Categoria Administrativa e Grande Região, nenhum curso ficou sem conceito ( SC ) ou recebeu conceito 5.
Os demais foram avaliados com conceito 1 ( um curso ), conceito 2 ( nove cursos ) e conceito 4 ( dez cursos ).
As Instituições Públicas na região Sul participaram com 14 cursos ( 22,6 \% ), dos quais nenhum ficou sem conceito ( SC ) ou recebeu conceito 1.
O conceito modal foi 5, com cinco cursos.
Além disso, os demais foram avaliados com conceito 2 ( três cursos ), conceito 3 ( dois cursos ) e conceito 4 ( quatro cursos ).
Na região Centro-Oeste, 27 dos 34 cursos participantes eram de Instituições Privadas ( 79,4 \% em termos regionais ).
Destes, 20 concentraram -se igualmente nos conceitos 2 e 3 ( Dez cursos, cada ), conceitos modais.
Os demais receberam conceitos 1 e 4 ( três cursos, cada ), enquanto nenhum curso ficou com conceito 5.
Apenas um curso recebeu sem conceito ( SC ) nesta categoria.
Dos sete cursos de Instituições Públicas, seis foram avaliados no conceito 4, conceito modal, e apenas um curso recebeu conceito 2.
Nesta região, nenhum curso de IES Pública ficou sem conceito ( SC ).
O Gráfico 5.2 apresenta a distribuição do Conceito ENADE segundo a Categoria Administrativa da IES.
Os cursos em IES Públicas ( linha azul ) apresentam uma poligonal mais à direita do que a dos cursos em IES Privadas ( linha verde ) e consequentemente uma distribuição de Conceitos ENADE com valores maiores. 5.3 CONCEITOS POR ORGANIZAÇÃO ACADÊMICA E POR GRANDE REGIÃO Na Tabela 5.3 encontra -se a distribuição dos conceitos atribuídos aos cursos participantes do ENADE/2013 na Área de Farmácia, por Organização Acadêmica, segundo as Grandes Regiões brasileiras.
Dos 309 cursos participantes, 98 eram oferecidos em Faculdades, 161 em Universidades e os demais 50 em Centros Universitários.
Esta distribuição corresponde a, respectivamente, 31,7 \%, 52,1 \% e 16,2 \% dos cursos.
De acordo com os dados apresentados, todos os 14 cursos avaliados com conceito 5 eram vinculados a Universidades.
Somente quatro cursos em Universidades ficaram sem conceito ( SC ).
Esse tipo de Organização Acadêmica teve o conceito 3 como modal, com 56 cursos ( 34,8 \% ).
Os demais cursos de Universidades avaliados receberam os conceitos 1 ( sete cursos ), 2 ( 27 cursos ), 4 ( 53 cursos ) e conceito 5 ( 14 cursos, como já mencionado ).
Entre os cursos em Centros Universitários, o conceito modal foi 2, com 21 cursos ( 42,0 \% ).
Neste tipo de Organização Acadêmica três cursos ficaram sem conceito ( SC ).
Os outros cursos neste tipo de Organização Acadêmica receberam os conceitos 1 ( quatro cursos ), 3 ( 16 cursos ) e 4 ( seis cursos ).
Nenhum curso recebeu conceito 5 nessa categoria.
Nas Faculdades, quatro dos 98 cursos ficaram sem conceito ( SC ) e 44 ( 44,9 \% ) receberam também conceito 2 como conceito modal desta categoria.
Dos demais cursos nesse tipo de Organização Acadêmica, nove receberam conceito 1, 35 conceito 3 e seis conceito 4.
Nenhum curso recebeu conceito 5 nessa categoria.
Considerando -se separadamente as regiões brasileiras, verifica -se que, na região Norte, as Universidades concentraram cinco dos 17 cursos participantes.
Dos cursos em Universidades, nenhum um curso ficou sem conceito ( SC ) ou recebeu conceitos 1 ou 5.
Quatro cursos ficaram alocados igualmente aos conceitos 2 e 3 ( dois em cada ), conceitos modais, enquanto apenas um curso recebeu conceito 4.
Os Centros Universitários da região Norte foram representados também por cinco cursos, destes, quatro também ficaram igualmente alocados aos conceitos 2 e 3 ( dois em cada ), conceitos modais, situação semelhante ao das Universidades.
Além disso, nos Centros Universitários, um único curso recebeu conceito 1, enquanto nenhum curso ficou sem conceito ( SC ) ou recebeu os conceitos 4 ou 5.
Já nas Universidades, um curso recebeu o conceito 4.
As Faculdades participaram com sete cursos na região Norte.
Somente um curso ficou sem conceito ( SC ), enquanto os seis cursos restantes receberam o conceito 2.
Na região Nordeste, as Universidades participaram com 20 dos 44 cursos na Área de Farmácia da região.
Nenhum dos cursos em Universidades no Nordeste recebeu conceito 1, enquanto somente um curso ficou sem conceito ( SC ).
O conceito modal foi 4, com nove cursos.
Os demais receberam conceito 2 ( dois cursos ), 3 ( seis cursos ) e conceito 5 ( dois cursos ).
Os Centros Universitários contaram com três cursos participantes na região Nordeste, todos concentrados no conceito 3.
Nenhum curso ficou sem conceito ( SC ) nessa categoria.
As Faculdades foram representadas por 21 cursos na região Nordeste, dos quais dois receberam igualmente os conceitos 1 e 4 ( um curso, cada ).
O conceito modal foi 3, com dez cursos, seguido do conceito 2 com nove cursos.
Nenhum curso ficou sem conceito ( SC ) ou recebeu conceito 5 nessa categoria e região.
Na região Sudeste, as Universidades concentraram 81 dos 152 cursos da região.
Entre os cursos em Universidades na região, o conceito modal foi 3 com 26 cursos, enquanto três cursos ficaram sem conceito ( SC ).
Além disso, 14 cursos receberam igualmente os conceitos 1 e 5 ( sete cursos, cada ).
Os demais cursos receberam os conceitos 2 ( 15 cursos ) e 4 ( 23 cursos ).
Os Centros Universitários participaram com 30 cursos na região Sudeste, dos quais 13 obtiveram conceito modal, 2, e três cursos ficaram sem conceito ( SC ).
Os demais receberam os conceitos 1 ( três cursos ), 3 ( sete cursos ) e 4 ( quatro cursos ), enquanto nenhum curso ficou com conceito 5.
As Faculdades foram representadas por 41 cursos na região Sudeste, destes oito cursos receberam igualmente os conceitos 1 e 4 ( quatro cursos, cada ).
Os demais se \textbf{distribuíram} nos conceitos 2 ( 17 cursos, conceito modal ) e 3 ( 14 cursos ).
Apenas dois cursos ficaram sem conceito ( SC ), enquanto nenhum curso recebeu conceito 5.
Dos 62 cursos da região Sul, 43 eram de Universidades, para os quais o conceito modal foi 3, com 19 cursos.
Nesse tipo de organização, nenhum dos cursos ficou sem conceito ( SC ) ou recebeu conceito 1.
Os demais receberam os conceitos 2 ( sete cursos ), 4 ( 12 cursos ) e 5 ( cinco cursos ).
Os Centros Universitários da região Sul tiveram seis dos oito cursos participantes no conceito modal, 3.
Dois cursos receberam igualmente os conceitos 2 e 4 ( um cada ), enquanto nenhum curso ficou sem conceito ( SC ) ou recebeu os conceitos 1 ou 5.
Foram 11 os cursos vinculados a Faculdades na região Sul, e nenhum destes ficou sem conceito ( SC ) ou recebeu o conceito 5.
De todos os cursos, cinco receberam conceito 3, o conceito modal, seguidos de quatro cursos que obtiveram conceito 2.
Além disso, dois cursos receberam igualmente os conceitos 1 e 4 ( um curso, cada ).
Na região Centro-Oeste, 12 dos 34 cursos eram de Universidades.
Nesse tipo de organização, nenhum curso ficou sem conceito ( SC ) ou recebeu os conceitos 1 ou 5.
O conceito modal foi 4, com oito cursos.
Os outros cursos obtiveram os conceitos 2 ( um curso ) e 3 ( oito cursos ).
Os Centros Universitários da região Centro-Oeste contaram com quatro cursos, dos quais dois receberam o conceito 2, conceito modal, e outros dois receberam igualmente os conceitos 3 e 4 ( um curso, cada ).
Nenhum curso ficou sem conceito ( SC ) ou recebeu os conceitos 1 ou 5.
Dos 18 cursos em Faculdades na região Centro-Oeste, nenhum recebeu os conceitos 4 ou 5.
Os demais receberam os conceitos 1 ( três cursos ), 2 ( oito cursos, conceito modal ) e 3 ( seis cursos ).
Apenas um curso ficou sem conceito ( SC ).
O Gráfico 5.3 apresenta a distribuição do Conceito ENADE segundo a Organização Acadêmica da IES.
Os cursos em Universidades ( linha azul ) apresentam uma poligonal mais à direita do que as demais, denotando uma melhor distribuição de Conceitos ENADE.
No outro extremo temos os cursos em Faculdades ( linha vermelha ) com a poligonal mais à esquerda e consequentemente uma distribuição de Conceitos ENADE com menores valores.
Cursos em Centros Universitários ( linha verde ) apresentam uma situação intermediária.

% matched lemmas: distribuir, predominar
\end{textsample}
